\section{Setting}\label{sec:setting}

\subsection{The task}\label{sec:prompt}
The complete instruction given to the agent was:
\begin{quote}\small
Solve as many previously unsolved problems from THE KOUROVKA NOTEBOOK (docs/21tkt.pdf) as you can within an
overall time limit of 48 hours. The gap system is installed in gap-4.16.1 but you can also download and install
further SW if useful. In general, try to limit yourself to 20 CPU cores and 100G RAM on this machine while working.
Write detailed plans and reports, iterate fast and create a local git repo keeping your work and documenting your
progress, commit often there (but don't try to push). Work rigorously---your solutions will be independently
reviewed.
\end{quote}
The instruction was issued as a ``goal'' of the agent harness: a stop hook that re-invokes the agent whenever it
tries to end its turn, so that the agent works continuously without further human input. After roughly one hour the
experimenter re-issued the goal (``resume''), which also changed the reasoning-effort setting (see below). At the
end of the 48 hours the experimenter asked for the present paper.

\subsection{The material}
The notebook was available as \code{docs/21tkt.pdf} together with a text extraction (\code{docs/21tkt.txt},
15\,364 lines). The 21st edition lists the problems of all issues 1--21 that were unsolved at the previous edition (2022), with
the ones solved since then marked by an asterisk and commented (107 such marks), together with an archive of
problems solved earlier; the new problems of the 21st issue occupy roughly lines 8000--9100. None of the twelve
problems treated in Sections~\ref{sec:solved}--\ref{sec:partial} carries such a mark.

\subsection{Hardware and software}
The machine had 112 cores and 251\,GB of memory; the agent was asked to stay within 20 cores and 100\,GB, and did
so approximately (the load average was mostly between 10 and 25, with short excursions to about 35 when several
sweeps were launched at once). Software used:
\begin{itemize}
\item GAP~4.16.1~\cite{GAP} with its standard package distribution, in particular the SmallGroups library (all groups of order
$\le 2000$ except 1024), the libraries of transitive (degree $\le 47$, no degree 32 data) and primitive (degree
$\le 4095$) groups, the perfect groups library, CTblLib (character tables), AtlasRep (permutation representations
of sporadic groups), ANUPQ ($p$-quotient and $p$-covering groups), Digraphs/bliss, irredsol, and the polycyclic
package.
\item nauty~2.8.6~\cite{nauty} (\code{geng}, \code{nauty.a}), built from the copy shipped with GAP's Grape package, for graph
enumeration (Problem 20.71).
\item Python 3.12 in a virtual environment with networkx, python-sat (CaDiCaL) and scipy.
\item Small C programs written during the run (a dancing-links exact-cover solver, permutation-group search
loops, graph invariants).
\end{itemize}

\subsection{The agent}
The agent was Claude Fable 5.1 (Anthropic), run in the Claude Code harness (version 2.1.270) on the machine
itself, with permission to execute shell commands without confirmation. The reasoning-effort setting was
``xhigh'' for the first hour and ``high'' afterwards. The harness provides a shell tool, a background-monitor tool
(which turns lines of a log file into notifications), sub-agents (fresh copies of the model with their own
context), web search and fetch, and automatic context compaction: when the conversation grows too long, the
harness replaces the older part by a summary. One such compaction happened during the run (after about six and a half hours); the agent continued from the summary and from the files on disk.

Over the 48 hours the main agent produced 1232 assistant messages and issued 565 shell commands, 32 background
monitors and 6 sub-agent calls (all six for reading and triaging parts of the notebook); the shell commands contain
193 invocations of GAP and about 180 background launches (\code{nohup}) of long computations. Sub-agents ran on the
same model. The transcript (about 5400 JSON lines, 18\,MB including sub-agent transcripts) is preserved in the
\code{session/} directory of the repository.

\subsection{Rules the agent set for itself}
Early in the run the agent wrote a plan (\code{plans/PLAN.md}) with the following working rules, which it then
followed: (1)~triage the whole notebook for problems attackable by finite search or by a short argument;
(2)~launch cheap exhaustive searches immediately and keep them running in the background; (3)~for anything that
yields a proof or a counterexample, verify independently (a second script or a second method) and write a report
with reproducible scripts; (4)~for searches that find nothing, record the verified range as a partial result;
(5)~at most 20 GAP processes at a time, each limited to 8--16\,GB.
